\section{Operations}
\label{sec:operations}

\begin{question}
  \textbf{What does it cost to run production training here?} How long do checkpoints take to save and
  load, and how long is training blocked? Is throughput stable over an hour, and are some GPUs or nodes
  slower? How long until training starts? How many tokens does a kWh buy?
\end{question}

\paragraph{Checkpoints (F1).} Save and load of Qwen3.8-27B (about 0.38\,TB \est{}) at 1, 13 and 52 nodes,
and of GLM-5, weights only (about 1.5\,TB \est{}) and full (about 10.4\,TB \est{}), at 13 and 52 nodes;
synchronous vs.\ asynchronous saves (time training is blocked vs.\ total write time); 3 repetitions at
different times of day. The nodes have no local disk, so an asynchronous save can only stage in host
memory before writing to the shared VAST filesystem. That filesystem serves the rest of Trillium, so the
sizes escalate with the administrators' agreement (Qwen 1 node, Qwen 52 nodes, GLM-5 weights, GLM-5 full),
and only one GLM-5 checkpoint is on disk at a time, deleted once its load is timed.

At the largest scale, a synchronous save takes \TriCkptQwenSave\,s for Qwen3.8-27B and
\TriCkptGlmFullSave\,s for the full GLM-5 state; an asynchronous save of GLM-5 blocks training for
\TriCkptGlmFullBlocked\,s. The highest write throughput is \TriCkptWriteGBs\,GB/s.

\begin{figure}[h]
  \centering\fig{checkpoint}
  \caption{Checkpoint time vs.\ number of nodes for Qwen3.8-27B (full state), GLM-5 weights only and GLM-5
    full state: synchronous save, asynchronous save (total and time training is blocked), and load.
    Right: filesystem throughput of saves and loads. \syntheticnote{checkpoint}}
  \label{fig:checkpoint}
\end{figure}

\placeholder{variation with time of day; write throughput against the filesystem's rated bandwidth.}

\paragraph{Stability (F2).} The best Qwen layout trains for at least one hour on 52 nodes. GPU clocks and
power from the node telemetry share the time axis with throughput, so a drop can be traced to a GPU.

\begin{figure}[h]
  \centering\fig{stability}
  \caption{Long run (52 nodes), vs.\ time: tokens/s per GPU per step with a 15-step median (top); SM
    clock (middle) and power (bottom) of the GPUs: gray band, range over all 416 GPUs; line, the median;
    the GPU with the lowest value is named. \syntheticnote{stability}}
  \label{fig:stability}
\end{figure}

\placeholder{throughput drift over the hour; throttling events and their cause.}

\paragraph{Stragglers (F2).} With data parallelism every rank waits for the slowest one, so per-rank step
times are equal and cannot point to a slow GPU. The slow GPU is the one that waits least: its activity is
highest. A node whose CPU clock falls also holds back every rank on it.

\begin{figure}[h]
  \centering\fig{stragglers}
  \caption{Long run, from the node telemetry. Top: activity of each GPU (highest 10-minute median), by
    node. Bottom: per node, the lowest 1-minute median of the CPU busy-core clock relative to the first
    minute. Nodes grouped by subnet; outliers are named. \syntheticnote{stragglers}}
  \label{fig:stragglers}
\end{figure}

\placeholder{slowest GPUs and nodes, and their effect on throughput.}

\paragraph{Startup (F3).} Time from process start to the first step, taken from every run (no extra
runs). At the largest scale it is \TriStartupQwen\,min for Qwen3.8-27B and \TriStartupGlm\,min for
GLM-5.

\begin{figure}[h]
  \centering\fig[0.72\linewidth]{startup}
  \caption{Time from process start to the first step vs.\ number of nodes, for both workloads.
    \syntheticnote{startup}}
  \label{fig:startup}
\end{figure}

\placeholder{which phase (environment, initialization, compilation, first step) grows with scale.}

\paragraph{Energy (F4).} Tokens per kWh, from the GPU power sampled by the node telemetry in every run (no
extra runs). Only GPU power is counted: CPUs, memory, network and cooling are not. The most efficient
Qwen layout at 13 nodes reaches \TriEnergyBest{} million tokens per kWh.

\begin{figure}[h]
  \centering\fig{energy}
  \caption{Million tokens per kWh (GPU power only): Qwen3.8-27B by layout at 13 nodes and 24k tokens per
    GPU (left); by workload and precision on the whole machine (middle); Qwen3.8-27B weak scaling vs.\
    nodes (right). \syntheticnote{energy}}
  \label{fig:energy}
\end{figure}

\placeholder{energy per token of the best configurations; effect of precision and scale.}
